%===================================================================================
% JORNADA CIENTÍFICA ESTUDIANTIL - MATCOM, UH
%===================================================================================
% Esta plantilla ha sido diseñada para ser usada en los artículos de la
% Jornada Científica Estudiantil de MatCom.
%
% Por favor, siga las instrucciones de esta plantilla y rellene en las secciones
% correspondientes.
%
% NOTA: Necesitará el archivo 'jcematcom.sty' en la misma carpeta donde esté este
%       archivo para poder utilizar esta plantila.
%===================================================================================



%===================================================================================
% PREÁMBULO
%-----------------------------------------------------------------------------------
\documentclass[a4paper,10pt]{article}

%===================================================================================
% Paquetes
%-----------------------------------------------------------------------------------
\usepackage{amsmath}
\usepackage{amsfonts}
\usepackage{amssymb}
\usepackage{jcematcom}
% Compatibility: neutralize tutor/title header behavior from the style so the
% template doesn't print "Tutor(es):" or the Jornada/tribunal heading.
% We redefine the internal macros to be no-ops here (local, safe change).
\makeatletter
% Don't print the tutor line in the title block
\@ifundefined{@starttutors}{\newcommand{\@starttutors}{}}{\renewcommand{\@starttutors}{}}
% Ensure there is no tutor content
\@ifundefined{@tutors}{\newcommand{\@tutors}{}}{\renewcommand{\@tutors}{}}
% Neutralize the automatic journal heading (avoid 'Jornada Científica...' text)
\@ifundefined{jcematcomheading}{\newcommand{\jcematcomheading}[3]{\thispagestyle{plain}}}{\renewcommand{\jcematcomheading}[3]{\thispagestyle{plain}}}
\makeatother
\providecommand{\tutors}{}
\usepackage[utf8]{inputenc}
\usepackage{listings}
\usepackage[pdftex]{hyperref}
\usepackage{caption}
\usepackage{subcaption}
\usepackage{graphicx}
\usepackage{booktabs}
\usepackage{float}
\usepackage{booktabs}
%-----------------------------------------------------------------------------------
% Configuración
%-----------------------------------------------------------------------------------
\hypersetup{colorlinks,%
	    citecolor=black,%
	    filecolor=black,%
	    linkcolor=black,%
	    urlcolor=blue}

%===================================================================================



%===================================================================================
% Presentacion
%-----------------------------------------------------------------------------------
% Título
%-----------------------------------------------------------------------------------
\title{Proyecto Conjunto de Ecuaciones Diferenciales Ordinarias y Matemática Numérica}

%-----------------------------------------------------------------------------------
% Autores
%-----------------------------------------------------------------------------------
\author{\\
\name Ernesto J. Govea Varona \email \href{mailto:ernesto.govea@estudiantes.matcom.uh.cu}{ernesto.govea@estudiantes.matcom.uh.cu}
	\\ \addr Grupo C211 \AND
\name Anthony Cruz García \email \href{mailto:anthony.cruz.estudiantes.matcom.uh.cu}{anthony.cruz.estudiantes.matcom.uh.cu}
	\\ \addr Grupo C211 \AND
\name Olivia Ortiz Arboláez \email \href{mailto:olivia.ortiz@estudiantes.matcom.uh.cu}{olivia.ortiz@estudiantes.matcom.uh.cu}
	\\ \addr Grupo C211
}
%-----------------------------------------------------------------------------------

%-----------------------------------------------------------------------------------

%-----------------------------------------------------------------------------------
% Headings (template heading neutralized above)
%-----------------------------------------------------------------------------------
% Original: \jcematcomheading{\the\year}{1-\pageref{end}}{A. Uno, A. Dos}
% The style's heading was neutralized; leave short headings blank so nothing
% template-specific appears in the page headers.
\ShortHeadings{}{}
%===================================================================================



%===================================================================================
% DOCUMENTO
%-----------------------------------------------------------------------------------
\begin{document}

%-----------------------------------------------------------------------------------
\maketitle

\vspace{-0.1cm}
% --- Equipo y tema (mostrar debajo del título)
\begin{center}
  {\large\bfseries Equipo --- Tema 4: Enfriamiento y calentamiento}
\end{center}
\vspace{1.0cm}

%===================================================================================
% Resumen y Abstract
%-----------------------------------------------------------------------------------
\selectlanguage{spanish} % Para producir el documento en Español

%-----------------------------------------------------------------------------------
% Resumen en Español
%-----------------------------------------------------------------------------------
\begin{abstract}
	Este proyecto aplica modelos basados en ecuaciones diferenciales ordinarias y
	métodos numéricos al estudio del enfriamiento post‑mortem y la estimación del
	tiempo de la muerte a partir de mediciones de temperatura. En la Parte A se
	utiliza la ley de Newton para ajustar datos y estimar parámetros, complementada
	con isoclinas y análisis del campo de pendientes. La Parte B analiza un
	modelo unidimensional no lineal dependiente del parámetro $\mu$: se identifican
	los equilibrios, se estudia su estabilidad y se construye el diagrama de
	bifurcación. La Parte C propone un modelo de dos compartimentos (núcleo y
	capa externa), realiza un análisis espectral para clasificar puntos críticos
	e interpreta las trayectorias en el plano de fase. El trabajo combina
	resultados analíticos y simulaciones numéricas (comparación de esquemas
	numéricos) y discute las implicaciones prácticas y limitaciones del enfoque
	en contextos forenses.
	
\end{abstract}

% English Abstract
\begin{enabstract}
	We present a concise study combining ordinary differential equation models
	and numerical methods to analyze post‑mortem cooling and support time‑of‑death
	estimation. Part A fits Newton's law to temperature measurements and
	interprets results using isoclines and slope fields. Part B studies a
	one‑dimensional nonlinear model parametrized by $\mu$, locating equilibria,
	assessing their stability and illustrating regime changes via a bifurcation
	diagram. Part C formulates a two‑compartment (core–shell) model, classifies
	critical points by eigenanalysis and interprets phase‑plane trajectories.
	Numerical comparisons (Euler, improved Euler, RK4) complement the analysis
	and the paper discusses practical implications and limitations for
	forensic applications.
\end{enabstract}
% Temas
%-----------------------------------------------------------------------------------
\begin{topics}
	Enfriamiento post-mortem y modelos de transferencia de calor.\\
	Modelos compartimentales; análisis de estabilidad (isoclinas y plano de fase) y métodos numéricos aplicados.
\end{topics}

\vspace{5.0mm}

%-----------------------------------------------------------------------------------
% NO BORRAR ESTAS LINEAS!
%-----------------------------------------------------------------------------------


%-----------------------------------------------------------------------------------

% --- Minipage con imágenes forense y clock (entre abstract y introducción)
\begin{center}
\begin{minipage}{0.28\textwidth}
	\centering
	\IfFileExists{Figures/forense.png}{%
		\includegraphics[width=0.90\linewidth]{Figures/forense.png}%
	}{%
		\small\texttt{Imagen Figures/forense.png no encontrada.}%
	}
	\vspace{-1.5mm}
	\small Figura: Forense
\end{minipage}\hspace{0.02\textwidth}
\begin{minipage}{0.28\textwidth}
	\centering
	\IfFileExists{Figures/clock.png}{%
		
\includegraphics[width=0.90\linewidth]{Figures/clock.png}%
	}{%
		\small\texttt{Imagen Figures/clock.png no encontrada.}%
	}
	\vspace{-1.5mm}
	\small Figura: Reloj
\end{minipage}
\end{center}

\newpage

%===================================================================================

%===================================================================================
% Introducción
%-----------------------------------------------------------------------------------
\section{Introducción}\label{sec:intro}
%-----------------------------------------------------------------------------------
	Los procesos de transferencia de calor intervienen en fenómenos de la vida
	diaria y en aplicaciones tecnológicas y forenses. Las ecuaciones diferenciales
	ordinarias (EDO) proporcionan un lenguaje compacto y eficaz para describir
	estas dinámicas: permiten obtener soluciones explícitas en modelos lineales y
	ofrecen, mediante herramientas cualitativas (campos de pendientes, isoclinas,
	planos de fase y análisis de bifurcaciones), una comprensión profunda cuando
	aparecen no linealidades o acoplamientos entre componentes.

	El propósito de este trabajo es aplicar y comparar técnicas analíticas y
	cualitativas de EDO en problemas térmicos sencillos pero representativos. A
	partir del enfoque clásico presentado en Edwards \& Penney (4.ª ed., págs.
	40--41) se estudian ejemplos que ilustran cómo extraer conclusiones físicas a
	partir de modelos matemáticos: estimación de parámetros a partir de datos,
	interpretación de cambios de régimen ligados a parámetros del sistema y
	análisis de la dinámica acoplada en modelos de múltiples compartimentos. La
	presentación se orienta a ser útil tanto en contextos docentes como en
	aplicaciones prácticas, priorizando la claridad y la interpretación física
	sobre la enumeración técnica de resultados.

%===================================================================================



%===================================================================================
% Desarrollo
%-----------------------------------------------------------------------------------
\section{Desarrollo}\label{sec:dev}
%-----------------------------------------------------------------------------------%-----------------------------------------------------------------------------------
	\subsection{Parte A --- Ley de Newton e isoclinas}\label{sub:parteA}

	\subsubsection{Planteamiento y datos}
	\medskip
	\begin{itemize}
		\item Temperatura ambiente: $T_a = 70\,^{\circ}\mathrm{F}$ (constante).
		\item Temperatura corporal inicial (momento de la muerte): $T(0)=98.6\,^{\circ}\mathrm{F}$.
		\item Mediciones: $T=80\,^{\circ}\mathrm{F}$ a las 12:00 y $T=75\,^{\circ}\mathrm{F}$ a la 13:00.
	\end{itemize}

	\subsubsection{Solución}
	\medskip
	Por la ley de Newton del enfriamiento:
	\begin{equation}\label{eq:newton}
		\frac{dT}{dt} = -k\,(T-T_a),\qquad k>0.
	\end{equation}

	La solución general (con $t$ en horas desde la muerte) es
	\begin{equation}\label{eq:solNewton}
		T(t)=T_a + (T(0)-T_a)e^{-k t} = 70 + 28.6\,e^{-k t}.
	\end{equation}

	Denotando por $t_1$ las horas transcurridas entre la muerte y las 12:00 se obtiene
	las dos ecuaciones siguientes para los datos medidos:
	\begin{align}
		T(t_1) &= 70+28.6\,e^{-k t_1}=80 \\%
		&\Rightarrow\quad e^{-k t_1}=\frac{10}{28.6}, \notag\\
		T(t_1+1) &= 70+28.6\,e^{-k(t_1+1)}=75 \\%
		&\Rightarrow\quad e^{-k(t_1+1)}=\frac{5}{28.6}. \notag
	\end{align}

	Dividiendo ambas expresiones se obtiene $e^{-k}=\tfrac{1}{2}$, por tanto
	\begin{equation}
		k = \ln 2 \approx 0.6931\ \mathrm{h}^{-1}.
	\end{equation}

	Sustituyendo en cualquiera de las expresiones para recuperar $t_1$:
	\begin{equation}
		t_1 = -\frac{\ln(10/28.6)}{\ln 2} \approx 1.515\ \mathrm{h} \approx 1\ \mathrm{h}\;30\ \mathrm{min}\;54\ \mathrm{s}.
	\end{equation}

	Por tanto, en el marco del modelo y las hipótesis formuladas, la estimación de la hora de la muerte es aproximadamente 1.515 horas antes de las 12:00 (10:29 AM).

	\medskip
	\paragraph{Interpretación cualitativa e isoclinas}
	La ecuación (\ref{eq:newton}) depende únicamente de $T$, por lo que las isoclinas en el plano $(t,T)$ son líneas horizontales. El campo de pendientes indica que:
	\begin{itemize}
		\item Si $T>T_a$ entonces $dT/dt<0$ (enfriamiento hacia $T_a$).
		\item Si $T<T_a$ entonces $dT/dt>0$ (calentamiento hacia $T_a$).
		\item En $T=T_a$ la derivada es nula; $T_a$ es un equilibrio estable.
	\end{itemize}

	Las soluciones decaen exponencialmente hacia $T_a$ sin cruzarse; el tiempo característico del proceso es $1/k$.
	\newpage
% --- Figura: isoclinas y campo de pendientes (Parte A)
\begin{figure}[htb]
	\centering
		\IfFileExists{Figures/isoclinas_partA.png}{%
			\includegraphics[width=0.95\textwidth]{Figures/isoclinas_partA.png}%
		}{%
			\centering\small\texttt{Imagen Figures/isoclinas\_partA.png no encontrada.}%
		}
	\caption{Isoclinas y campo de pendientes para la Ley de Newton (Parte A).}
	\label{fig:isoclinas_partA}
\end{figure}

\paragraph{Análisis Numérico}
Se compararon los métodos de Euler, Euler Mejorado y Runge-Kutta 4 (RK4) con la solución analítica $T(2)=72.5\,^{\circ}\mathrm{F}$.

\begin{table}[htb]
\centering
\caption{Comparación numérica en $t=2$ horas}
\begin{tabular}{l c r r l}
	\toprule
	Método & $N$ & $T(2)\,(^{\circ}\mathrm{F})$ & Error & Costo \\
	\midrule
	Analítica & N/A & 72.500000 & 0.000000 & N/A \\
	Euler & 1 & 72.516834 & 0.016834 & N \\
	Euler Mejorado & 2 & 72.500138 & 0.000138 & 2N \\
	Runge-Kutta 4 & 4 & 72.500000 & 0.000000 & 4N \\
	\bottomrule
\end{tabular}
\end{table}

RK4 ofrece la mayor precisión a un costo computacional eficiente.

 
\subsection{Parte B --- Bifurcación}\label{sub:parteB}
\medskip
\subsubsection{Resolución:} 
Considere la EDO adimensional,
\begin{equation}\label{eq:bif}
	u' = \mu\,u + u^{3},\qquad \mu\in\mathbb{R}.
\end{equation}

Defina $f(u;\mu)=\mu u + u^{3}$; entonces $f'(u)=\mu + 3u^{2}$. Los puntos de equilibrio se obtienen resolviendo
\begin{equation}
	\mu\,u + u^{3} = 0 \quad\Longrightarrow\quad u(\mu + u^{2})=0.
\end{equation}

Por tanto las soluciones estacionarias son:
\begin{itemize}
	\item $u=0$, válida para todo $\mu$.
	\item $\mu + u^{2}=0\ \Rightarrow\ u=\pm\sqrt{-\mu}$, que existen sólo para $\mu<0$.
\end{itemize}

\paragraph{Estabilidad (análisis lineal)}
Linearizando alrededor de un equilibrio $u^*$ se usa $f'(u^* )$:
\begin{itemize}
	\item En $u=0$: $f'(0)=\mu$. Luego $u=0$ es linealmente estable si $\mu<0$ y linealmente inestable si $\mu>0$. En $\mu=0$ la linealización falla y el término cúbico domina: para pequeñas perturbaciones $\dot u\approx u^{3}$, lo que muestra inestabilidad no lineal en $u=0$.
	\item En $u=\pm\sqrt{-\mu}$ (cuando $\mu<0$): usando $u^{2}=-\mu$ se obtiene
	\[f'\big(\pm\sqrt{-\mu}\big)=\mu+3(-\mu)=-2\mu>0,\]
	por lo que estas ramas no triviales son linealmente inestables.
\end{itemize}

\paragraph{Interpretación y diagrama}
La rama trivial $u=0$ existe para todo $\mu$ y cambia su estabilidad en $\mu=0$. Las ramas no triviales $\pm\sqrt{-\mu}$ aparecen para $\mu<0$ y son inestables; en $\mu=0$ dichas ramas colapsan en la trivial. Debido a la simetría $u\mapsto -u$ y la estructura observada, la bifurcación corresponde a un pitchfork subcrítico: al aumentar $\mu$ a través de $0$ la solución trivial pierde estabilidad y las ramas no triviales inestables se anulan en el punto de bifurcación.

% --- Figura: diagrama de bifurcación (insertada tras Parte B)
\begin{figure}[htb]
	\centering
	\IfFileExists{Figures/bifurcation_mu_u.png}{%
		\includegraphics[width=0.95\textwidth]{Figures/bifurcation_mu_u.png}%
	}{%
		\centering\small\texttt{Imagen Figures/bifurcation\_mu\_u.png no encontrada.}%
	}
	\caption{Diagrama de bifurcación para $\dot u=\mu u + u^{3}$. Las ramas
		no triviales $u=\pm\sqrt{-\mu}$ aparecen para $\mu<0$ y son inestables; la
		rama trivial $u=0$ cambia su estabilidad en $\mu=0$.}
	\label{fig:bifurcation}
\end{figure}

\subsection{Parte C --- Plano de fase y Estabilidad}\label{sub:parteC}

\subsubsection*{Problema}
Considere el siguiente sistema lineal numérico que modela un sólido con temperatura interna $X(t)$ en contacto con una ``capa'' o camisa $Y(t)$, estando el ambiente fijo en $A = 70$ (se trabaja con temperaturas relativas $x = X - A$, $y = Y - A$):

\begin{equation}
	\begin{cases}
		\dfrac{dx}{dt} = -4(x - y), \\
		\dfrac{dy}{dt} = -y
	\end{cases}
\end{equation}

\begin{enumerate}
	\item Calcule los puntos críticos y clasifíquelos (tipo y estabilidad).
	\item Verifíquelo construyendo el plano de fase y explique brevemente el significado físico de las trayectorias.
\end{enumerate}

\subsubsection*{Solución}

\subsubsection*{Puntos críticos}

Para encontrar los puntos críticos, igualamos a cero ambas ecuaciones:

\begin{align}
	\frac{dx}{dt} &= -4(x - y) = 0 \label{eq:dx} \\
	\frac{dy}{dt} &= -y = 0 \label{eq:dy}
\end{align}

De la ecuación (\ref{eq:dy}): $-y = 0 \implies y = 0$

Sustituyendo en (\ref{eq:dx}): $-4(x - 0) = 0 \implies -4x = 0 \implies x = 0$

Por lo tanto, el único punto crítico es:
\[
(x, y) = (0, 0)
\]

\subsubsection*{Clasificación del punto crítico}

Escribimos el sistema en forma matricial:
\[
\begin{pmatrix}
\dot{x} \\ \dot{y}
\end{pmatrix}
=
\begin{pmatrix}
-4 & 4 \\
0 & -1
\end{pmatrix}
\begin{pmatrix}
x \\ y
\end{pmatrix}
\]

Calculamos los autovalores de la matriz:
\[
A = \begin{pmatrix}
-4 & 4 \\
0 & -1
\end{pmatrix}
\]

El polinomio característico es:
\[
\det(A - \lambda I) = \det\begin{pmatrix}
-4-\lambda & 4 \\
0 & -1-\lambda
\end{pmatrix} = (-4-\lambda)(-1-\lambda) - (4)(0)
\]
\[
= (\lambda + 4)(\lambda + 1) = 0
\]

Los autovalores son:
\[
\lambda_1 = -4, \quad \lambda_2 = -1
\]

\begin{align}
	(\text{tr } A)^{2} &> 4 |A| \\
	25 &> 16
\end{align}

	\textbf{Clasificación:}
\begin{itemize}
	\item Ambos autovalores son reales y negativos
	\item $\lambda_1 = -4 < 0$, $\lambda_2 = -1 < 0$ , raíces de la ecuación característica.
	\item Tipo de equilibrio: \textbf{Nodo estable} , \[(tr A)^2 > 4 |A| \]
	\item Estabilidad: \textbf{Asintóticamente estable}
\end{itemize}

\subsubsection*{Significado físico}

El punto crítico $(0, 0)$ corresponde a:
\[
x = X - 70 = 0 \implies X = 70
\]
\[
y = Y - 70 = 0 \implies Y = 70
\]

Esto significa que el sistema alcanza el equilibrio térmico cuando tanto el sólido como la capa tienen la misma temperatura del ambiente ($70^\circ$). Las trayectorias representan cómo las temperaturas relativas $x$ y $y$ decaen exponencialmente hacia cero, indicando que el sistema tiende al equilibrio térmico con el ambiente.
% --- Figura: diagrama de fase (Parte C)

\begin{figure}[htb]
	\centering
	\IfFileExists{Figures/phase_diagram.png}{%
		\includegraphics[width=0.7\textwidth]{Figures/phase_diagram.png}%
	}{%
		\centering\small\texttt{Imagen Figures/phase\_diagram.png no encontrada.}%
	}
	\caption{Diagrama de fases del sistema lineal de la Parte C. El origen es un nodo estable (autovalores reales y negativos).}
	\label{fig:phase}
\end{figure}


%===================================================================================
% Conclusiones
%-----------------------------------------------------------------------------------
\section{Conclusiones}\label{sec:conc}

	A continuación se presentan las conclusiones principales derivadas de este estudio, formuladas como observaciones aplicables al modelado y análisis del enfriamiento post-mortem:

\begin{enumerate}
	\item Los modelos matemáticos basados en ecuaciones diferenciales son herramientas eficaces para describir la dinámica térmica post-mortem. Al relacionar parámetros como la constante de enfriamiento, las condiciones ambientales y el acoplamiento entre compartimentos, es posible predecir curvas de temperatura que sirven de base para estimar el tiempo de la muerte.

	\item La inclusión de efectos físicos relevantes —heterogeneidad térmica entre núcleo y capa externa, tasas de intercambio entre compartimentos y temperatura ambiental variable— permite construir modelos acoplados más realistas. Estas extensiones mejoran la capacidad predictiva y la interpretación de las trayectorias térmicas observadas en la práctica forense.

	\item El análisis cualitativo (isoclinas, planos de fase y estudio de autovalores) aporta información esencial sobre la estabilidad de los equilibrios térmicos y las escalas temporales de relajación (modos rápidos y lentos). Este entendimiento geométrico es útil para diseñar protocolos de muestreo y para explicar por qué ciertas trayectorias convergen más rápidamente al equilibrio.

	\item En el ámbito numérico, la comparación entre esquemas (Euler, Euler mejorado, Runge–Kutta) y la validación frente a soluciones analíticas o datos experimentales son indispensables. Métodos de orden superior, como RK4, ofrecen precisión superior con un coste computacional razonable, pero la elección del paso y del método debe ajustarse a la calidad y resolución de los datos disponibles.

	\item Aunque los modelos simples facilitan el análisis, tienen limitaciones: suposiciones como propiedades térmicas constantes, homogeneidad del cuerpo y condiciones ambientales invariables pueden introducir sesgos en la estimación del tiempo de muerte. Por ello es crítico estimar parámetros a partir de datos y considerar la incorporación de efectos adicionales (ropa, masa corporal, variación ambiental) para mejorar la aplicabilidad forense.
\end{enumerate}

%===================================================================================



%===================================================================================
% Recomendaciones
%-----------------------------------------------------------------------------------
\section{Recomendaciones}\label{sec:rec}

	A continuación se presentan recomendaciones concretas para continuar y ampliar
	el trabajo desarrollado:

	\begin{enumerate}
		\item Realizar campañas experimentales controladas para recoger series temporales
		de temperatura en cuerpos o análogos (phantoms) bajo distintas condiciones de
		masa, ropa y exposición ambiental. Estos datos permitirán ajustar y validar los
		parámetros del modelo y cuantificar la incertidumbre en las estimaciones del
		tiempo de muerte.

		\item Extender el modelo simple de dos compartimentos incorporando dependencias
		temporales de la temperatura ambiente y propiedades térmicas que dependan de la
		temperatura o del tiempo (p. ej. conductividad variable, radiación). Evaluar
		el impacto de estas extensiones sobre las predicciones.

		\item Integrar resultados en una guía práctica para peritos forenses, describiendo
		limitaciones del modelo, intervalos de confianza en las estimaciones y condiciones
		de aplicabilidad (p. ej. masa corporal, ropa, viento). Esto acercará la metodología
		a la práctica pericial.
	\end{enumerate}

%===================================================================================



%===================================================================================
% Bibliografía
%-----------------------------------------------------------------------------------
\begin{thebibliography}{99}
%-----------------------------------------------------------------------------------
	\bibitem{zill} Dennis G. Zill. \emph{Ecuaciones diferenciales}. (Spanish edition of
		"Differential Equations with Boundary-Value Problems"). Cengage Learning, 2009.

	\bibitem{edwards_penney} C. Henry Edwards, David E. Penney. \emph{Differential
		Equations and Boundary-Value Problems: Computing and Modeling}. 4th ed.,
		Prentice Hall, 2008.

	\bibitem{henssge} C. Henssge, N. Madea. "Estimación del intervalo postmortem: métodos
		y procedimientos". En: manuales y revisiones forenses sobre termometría forense y
		estimación del tiempo de muerte (varias ediciones). (Consultas y capítulos de
		referencia para metodología experimental y práctica pericial.)

	\bibitem{strogatz} Steven H. Strogatz. \emph{Nonlinear Dynamics and Chaos: With Applications
		to Physics, Biology, Chemistry and Engineering}. 2nd ed., Westview Press, 2015.

	\bibitem{burden} Richard L. Burden, J. Douglas Faires. \emph{Numerical Analysis}. 9th ed.,
		Brooks/Cole (Cengage), 2011.

	\bibitem{press} W. H. Press, S. A. Teukolsky, W. T. Vetterling, B. P. Flannery.
		\emph{Numerical Recipes: The Art of Scientific Computing}. 3rd ed., Cambridge
		University Press, 2007.

	\bibitem{saferstein} Richard Saferstein and Tiffany Roy. \emph{Forensic Science: From the Crime Scene to the Crime Lab}, 4th Edition (eBook).
		(Text covering forensic laboratory methods, crime-scene processing and temporal estimation techniques.)

%-----------------------------------------------------------------------------------
\end{thebibliography}

%-----------------------------------------------------------------------------------

\label{end}

\end{document}

%===================================================================================
